\section{The substrate}
\label{sec:substrate}

The engine under everything that follows is \system{}, a
\texttt{no\_std} Rust spiking-neural-network library with an
\texttt{i16} fixed-point hot path (LIF neurons, STDP synapses, CSR
sparse connectivity), published as \texttt{neuralos-snn} on
crates.io. Integer-only computation is not an aesthetic choice: it is
the regime ternary weights live in natively (an \texttt{i16} state
holds a $\{-1,0,+1\}$ alphabet without conversion), and float-free SNN
computation has been independently motivated for FPU-less edge silicon
(IEEE 2025, \emph{Full-Integer SNN Inference with RISC-V ISA}). The
library forbids allocation and unwraps on the hot path; every number
in this paper is produced by that integer path.

\subsection{Two discretization grids, two dead zones}
\label{sec:grids}

Membrane potential integrates in integer units on one of two grids
(\texttt{VoltageResolution}): millivolt (the historical default) and
centi-millivolt (dead zone $\approx 2\,\mu$A), added when the mV
grid's quantization proved load-bearing. Each grid has a
\emph{per-type} dead zone for input currents (mV: $\sim$200\,$\mu$A
for E-cells, $\sim$100\,$\mu$A for I-cells) below which a single
pulse is absorbed without membrane movement. The climb barrier is
asymmetric and pinned by test: at a 300\,$\mu$A drive the membrane
sticks at $-59$\,mV; a $+12\,\mu$A recurrent pulse ratchets it
$+1$\,mV; a $-12\,\mu$A pulse is absorbed. Small weights are real
physics on these grids, not noise: whether a $\pm 12\,\mu$A pulse
moves the membrane is exactly the difference between the two grids,
and the paper's substrate claims are grid-keyed throughout.

\subsection{The population}
\label{sec:population}

All experiments in this paper use a 512-neuron slice of a real model
tensor --- 409 excitatory, 103 inhibitory, 261{,}632 directed pairs
--- with pair-based STDP ($\gamma = 125$), a fixed inhibitory wall
($I_{\mathrm{INH}} = 600\,\mu$A), and per-class raw/absorbed/applied
plasticity counters (Sec.~\ref{sec:mechanism}). Seeds are pinned;
every run is deterministic, and every claimed double-run produced
byte-identical exports.

\subsection{Case study: the dead wire, found by falsification}
\label{sec:deadwire}

The substrate's central debugging story is worth telling in full,
because it is the paper's clearest demonstration of
pre-registration paying for itself \emph{inside} the infrastructure,
months before the adjudication ladder did the same for the science.

\paragraph{Two pre-registered predictions, two falsifications.}
To test whether imported weights influence firing, we swept drive
amplitude with three nets --- imported weights, census-shuffled
control, zero --- and pre-registered the prediction that their spike
trains diverge at 600\,$\mu$A. Result: \textbf{honest NO} --- the
three trains were identical at \emph{every} amplitude (train Hamming
0, rate-L1 0, across the full 9-amplitude grid). A second, finer-ruler
sweep on the new centi grid (pre-registered again: divergence expected
once mV truncation stopped absorbing pulses) falsified the rescue
hypothesis too: totals shifted, E-cells fired at amplitudes the mV
grid had silenced --- and the trains remained identical.

\paragraph{The canary.}
Two falsified predictions forced a code re-read in place of a third
hypothesis, and the re-read found the structural cause:
\texttt{step()}'s Phase 2 injected the recurrent current
(\texttt{weight/10}) into the postsynaptic accumulator \emph{after}
Phase 1's integration, and the next step's opening loop called the
accumulator clear \emph{before} Phase 1 --- every recurrent pulse was
born and destroyed without ever reaching \texttt{integrate\_and\_fire}.
The canary test
(\texttt{recurrent\_current\_is\_never\_integrated\_in\_step\_bug\_pinned})
pinned it empirically: a 2-neuron net, the presynaptic cell fires, the
postsynaptic membrane never leaves $-70$\,mV. The wire had been cut
since the original port; the mV dead zone (real, unit-proven) had been
a co-blocker upstream of the cut. An earlier audit's claim that
``the pulse integrates with one-step delay'' was re-reading the code
comment's intent, not the code's behavior.

\paragraph{The fix and the shared miss.}
The fix reorders the step: adaptation-current decay $\to$ integrate
(reads the previous step's pulses --- the one-step delay the design
always claimed) $\to$ clear-after-read $\to$ propagate $\to$
plasticity. Exact-value tests pin transmission on both grids
(centi: $-7{,}000 \to -6{,}988$ exactly one step after the presynaptic
spike; mV: $-70 \to -68$ with a strong weight; two presynaptic spikes
sum their quanta). The chastening datum: the library's 152-test suite
was green before and after --- \emph{no unit pin had ever exercised
live transmission through \texttt{step()}}. The tests pinned the
parts; nothing pinned the wire between them.

\paragraph{The mechanism reversal.}
Re-running the import experiment on the live wire reversed the
substrate's plasticity mechanism. Dead-wire era: correlated pairs
drifted \emph{negative} (same-step co-fire tie-break LTD,
anti-Hebbian, $\Delta\mathrm{SI} = -1.0000$) --- an artifact of a
regime in which weights could not transmit. Live wire: pre spikes
causally drive post spikes one step later $\Rightarrow$
pre-before-post $\Rightarrow$ LTP: intra-class mean $\Delta$
flipped from $-0.3133$ to $+0.1075$, discrimination magnitude
$|\Delta\mathrm{SI}| = 1.0000$ in \emph{both} eras. The sign was
never the selectivity claim; it is the mechanism label. The frozen
gate's directional criterion --- written during the dead-wire era ---
read the post-fix result as COLLAPSES and parked the pipeline
honestly; the amended criterion asserts the raw non-degenerate fields
(intra $|\mathrm{mean}\,\Delta| \ge 0.05$, flips $>0$, Hamming
$<0.50$, zero sign crossings) and prints the direction as the era's
mechanism label.

\subsection{Arrangement vs.\ census}
\label{sec:arrangement}

With the wire live, the import experiment's question got a real
answer. On the centi grid at 600\,$\mu$A: imported vs.\ its own
census-shuffle control diverges far more than imported vs.\ zero ---
$H(i,c) = 39{,}011$ vs.\ $H(i,z) = 23{,}907 \approx H(c,z) =
23{,}958$ (1.63$\times$): \emph{firing reads the arrangement of the
model's weights, not just their census}. The mV grid tells the
complementary story: rate/Hamming sensitivity to weights (Hamming
58{,}779 at 600\,$\mu$A; recruitment at 300\,$\mu$A where imported
fires 1.14\,Hz, control 0.99, zero 0.00 --- coherent same-group
bursts stack past the shrunken margin; a Gaussian-$\sigma$ analysis
would have missed that bursts are coherent) --- with single weight
pulses dead-zone-absorbed. Under the model's own in-vivo drive the
same inequality holds at 1.35--1.36$\times$
(Sec.~\ref{sec:invivo}). Boundary, stated: one slice, one model, one
seed --- a sensitivity result, not a population claim.

\subsection{The mechanism label, counted}
\label{sec:mechanism}

The substrate's realized plasticity mechanism is not what the naive
Hebbian reading of Sec.~\ref{sec:deadwire} suggests --- and the
refutation came from the substrate's own counters. Cumulative
per-class counters ($\mathrm{raw} = \sum$ pre-clamp deltas;
$\mathrm{absorbed} = \sum$ clamped-away remainder; applied by
difference) decomposed the live-wire run: raw intra drift
$-739{,}295$ (mean $-17.85$/syn --- same-step and post-leads LTD
events dominate the raw stream) $\cdot$ clamp-absorbed $-839{,}029$
$\cdot$ \emph{applied} $+99{,}734$ (mean $+2.41$/syn). The realized
mechanism is \textsc{pairing-selective, clamp-rectified}: which pairs
drift is timing-driven (only co-firing pairs pair at all), but the
\emph{direction} of realized drift is bounds-driven --- the E-class
zero-floor absorbs the net-negative raw drift and the applied residue
potentiates. The examples compute the label from the counters (three
cases: raw-LTP-dominant $=$ Hebbian-carried; raw-negative with
applied-positive $=$ clamp-rectified; applied-negative $=$
LTD-carried); this is the decomposition the in-vivo runs of
Sec.~\ref{sec:invivo} inherit, and Fig.~\ref{fig:mechanism}
(Sec.~\ref{sec:loop}) sketches its two rectification stages.
